\documentclass[10pt]{article}
\usepackage{lingmacros}
\usepackage{tree-dvips}
\usepackage{makecell}
\begin{document}
\pagenumbering{gobble}
\section*{Componenti chimici della materia vivente}
Alcune proteine esistono in stato di costante disordine non ripiegato come... .\\
Le proteine nel momento del ripiegamento possono andare incontro ad un ripiegamento incompleto che porta ad agglomerati, un esempio è anche la \textbf{proteina prionica} che può passare da alfa-elica a beta-foglietto, quest'ultima struttura è più compatta e causa la malattia della mucca pazza; una proteina prionica malata può "infettare" quelle sane.\\
I proteoglicani sono proteine complesse coniugate a glicosamminoglicani (gruppi prostetici), costituiscono il collagene. Accumuli di queste macromolecole causano mucopolisaccarosidosi, un fattore di rischio è la mancanza genetica di enzimi nei lisosomi volti all'idrolisi dei mucopolisaccaridi.\\ 
Il \textbf{liquido extracellulare} funge da sito di appoggio per le cellule, le divide in tessuti diversi, le ciba immagazzinando fattori di crescita e permette di inviare stimoli (pressori) alle cellule. \\
I \textbf{mediatori chimici} che sfruttano i lipidi sfruttano ad esempio la separazione di \underline{acido arachidonico} dalla membrana cellulare nel momento in cui entra in contatto con un segnale esterno come i raggi ultravioletti, questo porta alla produzione di \underline{prostaglandine} che attivano il processo infiammatorio. Una produzione eccessiva di prostaglandine può causare melanomi. \\
\section*{La cellula}
È possibile visualizzare i componenti degli organelli cellulari attraverso il microscopio crioelettronico (4nm), mentre quelli ottici permettono di evidenziare solo l'interno della cellula. Sfruttando anticorpi con flurotrofi si possono evidenziare specifiche proteine come quelle del nucleo.\\
\textbf{Procariote} "\underline{prima del nucleo}", organismi che compongono il microbiota e che possono provocare patologie in caso di produzione eccessiva di metaboliti e scompensi. \\
Nelle cellule eucariote troviamo barriere lipidiche che mantengono compatte e sezionate le varie componenti, il \underline{nucleolo} è l'unico che non è delimitato ma è contenuto all'interno del nucleo stesso.\\
All'interno delle cellule la \underline{densità delle macromolecole} è molto alta, solo il poro nucleare è composto da più di 450 proteine.\\
L'origine degli organuli è avvenuta per invaginazione della membrana plasmatica nel caso del RE e fagocitosi di un altro microrganismo per il mitocondrio.
\paragraph*{Nucleo}Organello eucariote protetto da una doppia membrana tappezzata da pori nucleari che permettono all'mRNA e altre molecole di passare. Per evidenziare le 2 lamine nucleari sono stati usati fluorometri con la ricostruzione di ogni anello.\\
All'interno si trova il \underline{nucleoplasma} (simile al citosol) in cui sono immersi nucleolo e  cromatina. L'eucromatina si condensa a formare l'eterocromatina per formare i cromosomi a ridosso della membrana e del nucleolo, l'informazione viene presa solo dall'eucromatina rilassata.\\
Il \textbf{nucleolo} si occupa della produzione di \underline{rRNA} per produrre la subunità maggiore e minore dei ribosomi grazie a proteine ribosomali esterne che entrano nel nucleo, una volta prodotte le subunità queste escono dal nucleo per unirsi.
\paragraph*{Ribosomi}Si possono trovare liberi nei citosol dove producono proteine usate da perossisomi, nucleo e mitocondri; mentre quelli legati alla parete del reticolo endoplasmatico producono proteine che vengono passate all'apparato del golgi dove vengono fatte maturare per essere inviate alla membrana, vescicole...
\paragraph*{Reticolo endoplasmatico}Si divide in liscio e rugoso, quest'ultimo si appoggia al nucleo e si chiama così perché presenta su di esso i ribosomi, al suo interno le proteine diventano glicoproteine grazie alla glicosilazione. Se la proteine non si ripiega correttamente non viene rilasciata con vescicole (generata per gemmazione) ma viene riconosciuta e degradata dal proteasoma. \\
Il RE liscio si occupa della sintesi di lipidi e soprattutto di steroidi, contribuisce al metabolismo del glucosio nel fegato, detossifica sostanze tossiche e immagazzina $Ca^2+$.
\paragraph*{Apparato del Golgi}Accoglie nelle cisterne cis le vescicole del RE e successivamente vengono modificate nelle cisterne mediali e trans per essere poi inviate da qui alla membrana o all'esterno. Le vescicole del golgi sono usate per unirsi a vescicole esterne formando lisosomi.
\paragraph*{Lisosomi}Ambiente acido (pH 5) per attivare le idrolasi acide che obbliga le proteine interne ad essere particolarmente glicosilate. Si formano da lisosoma primario (trans golgi) e fagosoma (interno o esterno).\\
Nel \underline{sistema immunitario} digeriscono agenti patogeni grazie ai fagosomi che si fondono, una volta fagocitati agenti esterni, con i lisosomi formando autofagolisosomi.
\paragraph*{Perossisomi}Si occupano di beta e alfa ossidazione di acidi grassi e della produzione di sfingolipidi. 
\paragraph*{Mitocondri}Usati nella produzione di energia. I mitocondri si dividono per scissione e se sono danneggiati vengono digeriti da fagolisosomi nel processo selettivo di mitofagia.
\paragraph*{Citoscheletro}È fatto da microtubuli, microfilamenti e filamenti intermedi.\\
I \textbf{microtubuli} sono formati nel centrosoma da alfa e beta-tubulina disposti in cilindretti, sono polarizzati in maniera dinamica per aggiungere tubulina da una parte e perderla dall'altra. Compongono il fuso mitotico e sono la base di appoggio di chinesine (movimento anterogrado) e dineine (movimento retrogrado).\\
I \textbf{microfilamenti} sono formati da actina monoglobulare e danno sostegno alle cellule grazie alla sua organizzazione in \underline{fibre da stress}, inoltre modificano la forma della cellula con lamellipodi e filopodi per aiutare il movimento.\\
I \textbf{filamenti intermedi} sono solo negli organismi pluricellulari e sono costituiti da molecole diverse a seconda della sua posizione nell'organismo (cheratine o neurofilamenti). Sono presenti sotto la membrana nucleare formando le 2 \underline{lamine nucleari} che sostengono le membrane e danno forma al nucleo. 
\paragraph*{Matrice extracellulare}
Dona la forma alle cellule causando la morte cellulare nel caso si staccassero da questa. Può indurre segnali meccanici provocando rilassamento o contrazione dei filamenti d'actina nella cellula.
\paragraph*{Membrana plasmatica} Modello a mosaico fluido composto da proteine, fosfolipidi anfipatici e glicolipidi.Ogni membrana in base alla sua locazione ha una composizione diversa, anche le lunghezze dei singoli acidi grassi modifica la composizione, più sono lunghi più aa saranno necessari per comporre proteine transmembrana. 
\section*{Trasporto}
Il trasporto di sostanze più grandi richiede l'uso di vescicole le quali sono prodotte grazie all'associazione di filamenti di clatrina, una proteina che si unisce tra simili disponendosi a nido d'ape nella formazione delle vescicole. La clatrina è richiamata sulla membrana da --- che si possono concentrare in un punto per ragioni differenti come un'elevata concentrazione di recettori nel lato esterno.
\section*{Sintesi proteica}
Sintesi proteica e espressione genetica vengono modificate attraverso modifiche chimiche sugli istoni che vengono metilati e acetilati.\\
Nel caso della metilazione viene repressa l'espressione di un determinato gene, con la duplicazione del DNA la metilazione torna al suo stato originale interessando solo l'epigenetica e il fenotipo, mentre l'ereditarietà riguarda le mutazioni e le modifiche alla base del DNA.\\
L'acetilazione degli istoni aumenta la sintesi proteica, mentre la metilazione aumenta la condensazione in etero(?)cromatina. La metilazione avviene maggiormente fuori dalle CpG Island (zone concentrate di legami citosina-guanina), nel caso di cellule tumorali la metilazione avviene nelle CpG Island producendo fattori pro-oncogeni e impedendo una corretta modulazione del gene.\\
La sintesi proteica è regolata in particolare da miRNA, siRNA e altri sRNA, i miRNA si trovano nel sangue, latte materno e liquido cefalorachidiano. La SRP-RNA è una sequenza segnale che si appaia alle proteine per trasportare il ribosoma al RE.
\subsection*{Trascrizione} A monte del gene si trovano gli elementi distali di controllo della trascrizione composte da silencer e enhancer, mentre nei pressi del gene da trascrivere si trova il promotore e la TATA box (riconosciuta dalla RNA-polimerasi). Una proteina d'attivazione si lega all'enhancer che si lega al mediatore grazie ad un piegamento del segmento, il quale permette di formare un complesso di pre-inizio.
\paragraph*{mRNA} Una volta formato si circolarizza a formare un complesso di ribosomi che traducono la stessa sequenza chiamata \underline{poliribosoma} grazie a delle proteine che riconoscono cappuccio e coda del mRNA. Gli introni non codificanti vengono rimossi dallo \textbf{splicing} grazie agli spliceosomi composti(?) da snRNA che si occupa di riconoscere le estremità dell'introne dove tagliare per formare un cappio che viene rimosso. Lo splicing alternativo modifica la composizione degli esoni del mRNA producendo proteine diverse.
\subsection*{Traduzione}
\paragraph*{tRNA} Il codone di stop è trasportato da un fattore di rilascio e non un tRNA idrolizzando il legame tra tRNA nel P e il polipeptide. Gli chaperone molecolari impediscono l'aggregazione tra catene polipeptidiche nel poliribosoma. 
\section*{Replicazione}
\section*{Morte cellulare}
Le \textbf{caspasi} sono pro-enzimi che quando vengono attivate per taglio proteolitico innescano l'apoptosi, possono svolgere funzione di iniziatrici o effettrici. Possono essere attivate da recettori (estrinseca), citocromo c rilasciato dal mitocondrio (intrinseca) o un danno al DNA (gene p53 che attiva il mitocondrio). Il citocromo c permette di formare l'apoptosoma, mentre p53 è quel gene che previene la formazione di tumori nel caso di mutazioni (viene regolata dal MDM2, prodotto dal p53 stesso).\\
La \underline{necroptosi} è la morte cellulare programmata da cascata metabolica che sversa materiale cellulare all'esterno. Lo stesso segnale può indurre apoptosi o necroptosi. 
\section*{Stem cell e differenziamento}
I fattori di Yamanaka sono fattori retrovirali inducibili tramite retrovirus che permettono a una cellula di diventare pluripotente (IPSC).
\section*{Ciclo cellulare}
Lo zigote forma un organismo attraverso la divisione tra cellule staminali e cellule differenziate. \\
\textbf{Nei batteri} avviene per scissione binaria dove il DNA circolare viene duplicato e distribuito in egual misura nei due nuovi organismi.
\subsection*{Eucarioti}
Avviene con l'interfase (G1,S e G2) e fase M. Le cellule che non si riproducono differenziate (neuroni e cellule muscolari scheletriche) sono in G0 perenne, altri come fibroblasti e epatociti passano da G0 a G1 in seguito a uno stimolo, oltre a questi anche le cellule satellite staminali si attivano dopo un danno al muscolo per ristabilire le condizioni fisiologiche.
\paragraph*{G1} È la fase più lunga. La cellula decide se proliferare o entrare in fase G0, nel caso decidesse di riprodursi inizierebbe a accrescere e a duplicare i suoi organelli fino ad entrare in fase S quando è abbastanza grande. In questa fase la cromatina non è condensata, permettendo modifiche post-traduzionali che vengono sfruttate in fasi successive.
\paragraph*{S} Duplicazione del DNA in due cromatidi fratelli uniti dal centromero.
\paragraph*{G2} Aumento della sintesi proteica e duplicazione dei centrioli.
\paragraph*{Profase} Prima fase della mitosi, il cromosoma diventa più compatto e viene a formarsi il fuso mitotico.
\paragraph*{Prometafase} Inizia la disgregazione dell'involucro nucleare (anche grazie ai microtubuli) permettendo ai microtubuli di entrare in contatto con il cinetocore. Le proteine della membrana vengono disseminate nel RE. 
\paragraph*{Metafase} I cromatidi si dispongono sulla \textbf{piastra metafasica}, in questa fase viene analizzato il cariotipo perché i cromotidi sono compatti e condensati abbastanza da essere visibili. Vengono rimosse le coesine che tengono legate le due parti di cromatidi.
\paragraph*{Anafase} I cromosomi figli vengono spostai alle estremità della cellula dai microtubuli che si accorciano (depolimerizzazione all'estremitò positiva), il cinetocore viene mantenuto vicino all'estremità dei microtubuli grazie ad un complesso proteico che funge da ponte.
\paragraph*{Telofase} Ultima fase della mitosi. Iniziano a riorganizzarsi gli organelli e a ricomporsi la membrana nucleare (il RE entra in contatto con la cromatina e i tubuli compongono la membrana), all'interno del nucleo la cromatina si decondensa. 
\paragraph*{Citodieresi} Si forma un anello di acto-miosina contrattile che divide la membrana plasmatica.\\\\
I mitocondri si duplicano per scissione binaria grazie alla polimerasi gamma.\\
La divisione non è sempre simmetrica, il midbody ad esempio è l'ultima parte che mantiene unite le cellule figlie, quando avviene la separazione viene rilasciato oppure viene preso in maniera asimmetrica da una delle due figlie. La \textbf{divisione asimmetrica} è anche intenzionale come nel caso dell'oovogenesi e delle staminali (una staminale e una differenziata). Viene usata anche per rinnovare la cellula che forma una cellula principale e una più piccola contenente solo proteine danneggiate che rimangono nell'aggresoma, struttura piena di "rifiuti" che viene tenuta dalla cellula madre. È per questo motivo che le malattie causate da aggregati di proteine (Huntington, Parkinson, AD) colpiscono i neuroni, perché non possono rinnovarsi e formare una cellula sana.
\paragraph*{Checkpoint} Sono punti di controllo per assicurarsi che la cellula sia in grado di replicare: la prima è G1-S che si accerta sia abbastanza grande, poi G2-M si assicura che il DNA sia stato replicato e infine metafase-anafase controlla che i cromatidi siano in linea sulla piastra.\\
Questi controlli sono svolti dalla ciclina, una proteina sintetizzata che attiva la protein-chinasi ciclina dipendente che fosforila una proteina diversa per ogni fase approvata. Ogni checkpoint ha una ciclina diversa.\\
La ciclina e il suo segnale vengono poi terminati da una poli-ubiquitinazione.\\
Citochine e fattori di crescita stimolano la mitosi. I tumori inibiscono i checkpoint permettendo alla cellula di replicarsi all'impazzata, mentre soppressori come il p53 impediscono che questo avvenga.\\
\section*{Riproduzione}
\section*{Genetica}
Epistasi: due alleli su un locus regolano l'espressione di altri alleli, diversamente dall'esempio del gallo dove alleli e loci diversi interagiscono tra loro per dare un fenotipo intermedio.
\end{document}
